% GENERATED FILE. Edit canonical lessons, then run npm run generate:books.
% canonical-source-hash: fnv1a64:1cf68652cab23819
% canonical-lessons: SA-C472-talu, SA-C472-lala, SA-C472-carvayati, SA-C472-nisthivati, SA-C472-rasah

\chapter{The Mouth}
\label{ch:sa-the-mouth}
\hlchaptermodality{472}

\begin{chapteropening}
\textbf{By the end of this chapter:} \emph{I can talk about chewing and the mouth.}

Complete the last lesson of chapter 472: I can talk about chewing and the mouth.
\end{chapteropening}

\section[\texorpdfstring{\sk{तालु} (tālu)}{तालु (tālu)}]{\sk{तालु} (tālu) — the palate}

\label{lesson:SA-C472-talu}



\begin{quote}
Before the new one: say the Sanskrit for an eyelash, then the Sanskrit for shaves.
\end{quote}

\subsection*{You'll want to know: \sk{तालु}}
\textbf{\sk{तालु}} — \emph{tālu} — \textquotedblleft{}the palate\textquotedblright{}.

\begin{quote}
\sk{जिह्वा} \sk{तालु} \sk{स्पृशति।}
\end{quote}

\begin{quote}
\emph{jihvā tālu spṛśati.}
\end{quote}

\begin{quote}
The tongue touches the palate.
\end{quote}

A neuter noun.

\begin{cousinweb}
\textbf{\sk{तालु}} (\emph{tālu}) MW gives \textquotedblleft{}the palate\textquotedblright{}, from the Vājasaneyi Saṃhitā on; MW gives no derivation, so its origin is uncertain (MW).
\end{cousinweb}

\begin{grammarlens}[title={What you've built}]
One more word for describing what you lived and felt.
\end{grammarlens}

\subsection*{Your turn}
\begin{itemize}
\raggedright
  \item \emph{Say it:} \emph{tālu}
  \item \emph{Say it:} \emph{tālu}, once more
  \item \emph{Say it:} \emph{jihvā tālu spṛśati.} — The tongue touches the palate.
  \item \emph{Recall:} say the Sanskrit for an eyelash, then the Sanskrit for shaves, then say \emph{tālu} again
\end{itemize}

\begin{tcolorbox}[breakable,colback=teal!4,colframe=teal!35!black,title={Before you move on}]
What does \sk{तालु} mean? (\textquotedblleft{}The palate\textquotedblright{}.) Say it once more.
\end{tcolorbox}
\section[\texorpdfstring{\sk{लाला} (lālā)}{लाला (lālā)}]{\sk{लाला} (lālā) — saliva}

\label{lesson:SA-C472-lala}



\begin{quote}
Before the new one: say the Sanskrit for shaves, then the Sanskrit for the palate.
\end{quote}

\subsection*{You'll want to know: \sk{लाला}}
\textbf{\sk{लाला}} — \emph{lālā} — \textquotedblleft{}saliva\textquotedblright{}.

\begin{quote}
\sk{शुनकस्य} \sk{मुखात्} \sk{लाला} \sk{पतति।}
\end{quote}

\begin{quote}
\emph{śunakasya mukhāt lālā patati.}
\end{quote}

\begin{quote}
Saliva drips from the dog's mouth.
\end{quote}

A feminine noun.

\begin{cousinweb}
\textbf{\sk{लाला}} (\emph{lālā}) MW gives \textquotedblleft{}saliva, spittle, slobber\textquotedblright{}; MW gives no derivation, so its origin is uncertain (MW).
\end{cousinweb}

\begin{grammarlens}[title={What you've built}]
One more word for describing what you lived and felt.
\end{grammarlens}

\subsection*{Your turn}
\begin{itemize}
\raggedright
  \item \emph{Say it:} \emph{lālā}
  \item \emph{Say it:} \emph{lālā}, once more
  \item \emph{Say it:} \emph{śunakasya mukhāt lālā patati.} — Saliva drips from the dog's mouth.
  \item \emph{Recall:} say the Sanskrit for shaves, then the Sanskrit for the palate, then say \emph{lālā} again
\end{itemize}

\begin{tcolorbox}[breakable,colback=teal!4,colframe=teal!35!black,title={Before you move on}]
What does \sk{लाला} mean? (\textquotedblleft{}Saliva\textquotedblright{}.) Say it once more.
\end{tcolorbox}
\section[\texorpdfstring{\sk{चर्वयति} (carvayati)}{चर्वयति (carvayati)}]{\sk{चर्वयति} (carvayati) — chews}

\label{lesson:SA-C472-carvayati}



\begin{quote}
Before the new one: say the Sanskrit for the palate, then the Sanskrit for saliva.
\end{quote}

\subsection*{You'll want to know: \sk{चर्वयति}}
\textbf{\sk{चर्वयति}} — \emph{carvayati} — \textquotedblleft{}chews\textquotedblright{}.

\begin{quote}
\sk{बालः} \sk{शनैः} \sk{भोजनं} \sk{चर्वयति।}
\end{quote}

\begin{quote}
\emph{bālaḥ śanaiḥ bhojanaṃ carvayati.}
\end{quote}

\begin{quote}
The child chews his food slowly.
\end{quote}

Past (he …): \sk{सः} \sk{भोजनं} \sk{चर्वितवान्}, \emph{saḥ bhojanaṃ carvitavān}.

\begin{cousinweb}
\textbf{\sk{चर्वयति}} (\emph{carvayati}) comes from the root √carv, \textquotedblleft{}to grind with the teeth, to chew\textquotedblright{} (MW).
\end{cousinweb}

\begin{grammarlens}[title={What you've built}]
One more word for describing what you lived and felt.
\end{grammarlens}

\subsection*{Your turn}
\begin{itemize}
\raggedright
  \item \emph{Say it:} \emph{carvayati}
  \item \emph{Say it:} \emph{carvayati}, once more
  \item \emph{Say it:} \emph{bālaḥ śanaiḥ bhojanaṃ carvayati.} — The child chews his food slowly.
  \item \emph{Recall:} say the Sanskrit for the palate, then the Sanskrit for saliva, then say \emph{carvayati} again
\end{itemize}

\begin{tcolorbox}[breakable,colback=teal!4,colframe=teal!35!black,title={Before you move on}]
What does \sk{चर्वयति} mean? (\textquotedblleft{}Chews\textquotedblright{}.) Say it once more.
\end{tcolorbox}
\section[\texorpdfstring{\sk{निष्ठीवति} (niṣṭhīvati)}{निष्ठीवति (niṣṭhīvati)}]{\sk{निष्ठीवति} (niṣṭhīvati) — spits}

\label{lesson:SA-C472-nisthivati}



\begin{quote}
Before the new one: say the Sanskrit for saliva, then the Sanskrit for chews.
\end{quote}

\subsection*{You'll want to know: \sk{निष्ठीवति}}
\textbf{\sk{निष्ठीवति}} — \emph{niṣṭhīvati} — \textquotedblleft{}spits\textquotedblright{}.

\begin{quote}
\sk{मार्गे} \sk{न} \sk{निष्ठीवेत्।}
\end{quote}

\begin{quote}
\emph{mārge na niṣṭhīvet.}
\end{quote}

\begin{quote}
One should not spit on the road.
\end{quote}

Past (he …): \sk{सः} \sk{मार्गे} \sk{निष्ठ्यूतवान्}, \emph{saḥ mārge niṣṭhyūtavān}.

\begin{cousinweb}
\textbf{\sk{निष्ठीवति}} (\emph{niṣṭhīvati}) is built from \emph{ni}, \textquotedblleft{}down\textquotedblright{}, and the root √ṣṭhiv, \textquotedblleft{}to spit\textquotedblright{}; MW's ppp \emph{niṣṭhyūta} is \textquotedblleft{}spit out\textquotedblright{} (MW).
\end{cousinweb}

\begin{grammarlens}[title={What you've built}]
One more word for describing what you lived and felt.
\end{grammarlens}

\subsection*{Your turn}
\begin{itemize}
\raggedright
  \item \emph{Say it:} \emph{niṣṭhīvati}
  \item \emph{Say it:} \emph{niṣṭhīvati}, once more
  \item \emph{Say it:} \emph{mārge na niṣṭhīvet.} — One should not spit on the road.
  \item \emph{Recall:} say the Sanskrit for saliva, then the Sanskrit for chews, then say \emph{niṣṭhīvati} again
\end{itemize}

\begin{tcolorbox}[breakable,colback=teal!4,colframe=teal!35!black,title={Before you move on}]
What does \sk{निष्ठीवति} mean? (\textquotedblleft{}Spits\textquotedblright{}.) Say it once more.
\end{tcolorbox}
\section[\texorpdfstring{\sk{रसः} (rasaḥ)}{रसः (rasaḥ)}]{\sk{रसः} (rasaḥ) — a taste; a flavour}

\label{lesson:SA-C472-rasah}



\begin{quote}
Before the new one: say the Sanskrit for chews, then the Sanskrit for spits.
\end{quote}

\subsection*{You'll want to know: \sk{रसः}}
\textbf{\sk{रसः}} — \emph{rasaḥ} — \textquotedblleft{}a taste; a flavour\textquotedblright{}.

\begin{quote}
\sk{आम्रस्य} \sk{रसः} \sk{मधुरः} \sk{अस्ति।}
\end{quote}

\begin{quote}
\emph{āmrasya rasaḥ madhuraḥ asti.}
\end{quote}

\begin{quote}
The mango tastes sweet.
\end{quote}

A masculine noun.

\begin{cousinweb}
\textbf{\sk{रसः}} (\emph{rasaḥ}) MW's first sense is \textquotedblleft{}the sap or juice of plants\textquotedblright{}, the word in \textbf{\sk{फलरसः}}, and \textquotedblleft{}taste, flavour\textquotedblright{} comes later; MW gives no derivation, so its origin is uncertain (MW).
\end{cousinweb}

\begin{grammarlens}[title={What you've built}]
One more word for describing what you lived and felt.
\end{grammarlens}

\subsection*{Your turn}
\begin{itemize}
\raggedright
  \item \emph{Say it:} \emph{rasaḥ}
  \item \emph{Say it:} \emph{rasaḥ}, once more
  \item \emph{Say it:} \emph{āmrasya rasaḥ madhuraḥ asti.} — The mango tastes sweet.
  \item \emph{Recall:} say the Sanskrit for chews, then the Sanskrit for spits, then say \emph{rasaḥ} again
\end{itemize}

\begin{tcolorbox}[breakable,colback=teal!4,colframe=teal!35!black,title={Before you move on}]
What does \sk{रसः} mean? (\textquotedblleft{}A taste; a flavour\textquotedblright{}.) Say it once more.
\end{tcolorbox}
